% This file is part of Oaklisp.
%
% This program is free software; you can redistribute it and/or modify
% it under the terms of the GNU General Public License as published by
% the Free Software Foundation; either version 2 of the License, or
% (at your option) any later version.
%
% This program is distributed in the hope that it will be useful,
% but WITHOUT ANY WARRANTY; without even the implied warranty of
% MERCHANTABILITY or FITNESS FOR A PARTICULAR PURPOSE.  See the
% GNU General Public License for more details.
%
% The GNU GPL is available at http://www.gnu.org/licenses/gpl.html
% or from the Free Software Foundation, 59 Temple Place - Suite 330,
% Boston, MA 02111-1307, USA


\chapter{Miscellaneous} \label{misc}

\section{Tables}

The types are \df{generic-hash-table} and \df{eq-hash-table}.  The
access interface is \df{present?}, which returns a \df{pair} whose
\df{car} is the key and whose \df{cdr} is the associated value.  A
different interface to hash tables is provided by the T-style
\df{table-entry} operation which returns the associated value or
\df{\#f} if the key isn't in the table.  The setter of either
operation can be used to add, modify, and remove associations.

\makin{generic-hash-table}{key-hash-op equal-op}
\makin{eq-hash-table}{}

\so{present?}{table key}
\doc{Returns \texttt{(\emph{key \dt val})} pair, or \df{\#f} if not present.}

\so{table-entry}{table key}
\doc{Returns value indexed by \emph{key} or \df{\#f} if not present.}




\makin{equal-hash-table}{}
\doc{\df{generic-hash-table} with \df{equal?} for the comparison, so
that lists and strings that look alike are the same key.}

\fn{make-eq-hash-table}{}
\fn{make-equal-hash-table}{}
\fn{make-string-hash-table}{}
\doc{Return a new table of each kind.  A string table hashes the
characters, so two strings with the same text land in the same
bucket, and compares with \df{equal?}.}

\ty{hash-table}
\doc{The supertype of all of these.}

\section{Symbols}

\ty{symbol}
\pr{symbol?}{object}
\so{intern}{string}
\doc{The symbol with that name, made if there is none.  Symbols are
interned, so this is the only way two of them become \df{eq?}.}
\fn{gensym}{prefix}
\doc{A fresh symbol whose name is \emph{prefix} and a number.  It is
interned like any other, so a second \df{gensym} with the same prefix
is still a different symbol only because the number differs.}
\fn{generate-symbol}{prefix}
\fn{genvar}{}
\doc{\df{gensym} with the prefix \texttt{v}, which is what macro
expanders use for a temporary.}
\fn{genvar/1}{ignored}
\doc{\df{genvar} taking an argument, so it can be mapped over a list
to make one new name per element.}

\section{Characters}

\ty{character}
\pr{char?}{object}
\pr{graphic?}{character}
\doc{Whether the character has a printing glyph.}
\pr{digit?}{character radix}
\doc{Whether the character is a digit in that radix.}
\fn{digit-value}{character}
\doc{The number the character stands for as a digit, counting
\texttt{a} as 10 and upward; the character is not checked.}
\fn{digit->char}{n}
\doc{The other way, using upper case letters past 9.}
\op{upcase}{object}
\op{downcase}{object}
\doc{Defined on characters and on strings.}

\section{Delays}

Oaklisp's delays are compatible with the facility defined in R3RS, but
extend those primitive facilities in two ways.  First, the system will
automatically force promises when appropriate.  For instance, \texttt{(+
2 (delay 3))} does not signal an error; it returns \texttt{5}.
Similarly, delays are printed transparently, slightly violating
read/print consistency.  Secondly, the delay facility is user
extensible.  Users can create new kinds of delays that have special
protocols, for instance numeric delays that do not force themselves
upon arithmetic operations, but instead make more and more complicated
delays.


\sform{delay}{expression}
\doc{This immediately returns a \emph{promise} for \emph{expression},
without actually computing \emph{expression}.  This promise does not
compute \emph{expression} until it is forced, at which point it
returns the value of \emph{expression}, computing it if it hasn't
already done so.}

\op{force}{x}
\doc{If \emph{x} is a promise, it is forced to compute its value, which
is returned.  If \emph{x} is not a promise, it itself is returned.}

\ty{promise}
\doc{This is the type of the objects returned by \df{delay}.}

\ty{forcible}
\doc{This is an abstract type, of which \df{promise} is a concrete
subtype.  Subtypes of \df{forcible} are expected to respond to the
\df{force} operation in a sensible fashion.  Oaklisp's system
internals sometimes force instances of \df{forcible} automatically,
for instance when sending them messages for which no appropriate
method can otherwise be found.}

\fv{forcible-print-magic}
\doc{Controls how delays are printed.  This is how \texttt{(delay 'foo)}
would print under various settings of \texttt{forcible-print-magic}.
\begin{center}
\begin{tabular}{ll}\toprule
\textbf{value}             & \textbf{print style} \\ \midrule
\texttt{\#f}          & \texttt{\#<DELAY 3462>}   \\
\texttt{indicate}    & \texttt{\#[DELAY FOO 3462]}\\
\texttt{transparent} & \texttt{FOO}               \\ \bottomrule
\end{tabular}
\end{center}
The default is \df{transparent}.  A setting of \df{indicate} is more
instructive if you encounter odd behavior that might be due to
delays.}

\section{Threads}

Oaklisp's processes are green threads: they run in the one emulator
and switch at the points named here, and at alarms when the system
was built with preemption.  A process is a continuation and a set of
fluid bindings, so \df{bind} in one is not seen in another.

\ty{process}
\doc{A process descriptor.}
\fn{current-process}{}
\doc{The descriptor of the process this code is running in.}
\op{process-id}{process}
\fn{lwp}{thunk}
\doc{Queues \emph{thunk} to be run in the current process, which is
to say as a coroutine of it, sharing its fluid bindings.}
\fn{process-run-fn}{operation arguments}
\doc{Queues a call of \emph{operation} on \emph{arguments} in a
process of its own, with its own fluid bindings.}
\fn{pause}{}
\doc{Lets any other runnable process go, and returns when its turn
comes round again.}

\makin{mutex}{}
\op{acquire-mutex}{mutex}
\doc{Waits until the mutex is free and takes it.  This is a spin
lock, so a semaphore is better where the wait may be long.}
\op{release-mutex}{mutex}

\makin{test-and-settable}{}
\op{test-and-set}{test-and-settable predicate new-value}
\doc{Under the object's own mutex, applies \emph{predicate} to the
value it holds; if the answer is true the value becomes
\emph{new-value}, which is returned, and otherwise \df{\#f} is
returned.  Unlike the usual test-and-set the test is any predicate
and the new value any object, which is why a failure is told from a
success by the \df{\#f}.}

\makin{semaphore}{}
\op{wait}{semaphore}
\doc{Decrements the count, and blocks until a \df{signal} if it goes
negative.  A blocked process leaves the run queue, so this is the one
to wait on.}
\op{signal}{semaphore}
\doc{Increments the count and wakes the process that has waited
longest, if any.}

\section{The Outside World}

\op{get-time}{}
\doc{The CPU time used so far, in milliseconds.}
\mc{time}{\lpar n\rpar \dt body}
\doc{Runs \emph{body} \emph{n} times and returns the average time in
milliseconds.}

\fn{deep-describe}{object}
\doc{Like \df{describe}, but goes on to describe the contents of each
instance variable rather than just naming them.}

\gv{argline}
\doc{The command line arguments after the \texttt{--}, as a list of
strings.  This is set at startup from \df{get-argline}.}
\fn{get-argline}{}
\doc{Reads the argument list out of the emulator afresh.}
\so{cwd}{}
\doc{The current directory, as a string.  Setting it changes the
process's directory, and signals \df{error-changing-directory} if it
cannot.  The value is what Oaklisp last set, not what the operating
system would say, since there is no primitive that asks.}

\fn{fact}{n}
\doc{Factorial, kept around because it is what everyone types first.}
